\subsection{A matching logarithmic law for physical dictionaries}
\label{sec:r76-physical-approximation}

The preceding uppers use genuine density matrices as hidden dictionary
states.  Within that representation class a matching lower, including
the leading logarithmic constant, is possible.  This restriction is
substantive: an arbitrary classical hidden-state model can have terminal
decoder tables that are not Born probabilities of fixed density matrices.
The conclusion below does not impose that restriction on such a model
without proof.

Fix one commanded angle $\theta=2\pi\omega$ in the undamped instrument.
A \emph{physical-dictionary simulator} has finitely many labels with
fixed qubit density matrices $\sigma_i$, arbitrary stochastic report
transition tables, and terminal decoder
$d_E(i)=\tr(E\sigma_i)$ for every requested effect $E$.
The causal transition tables are fixed independently of any later
terminal-measurement request.
The transitions themselves need not implement the quantum instrument on
each dictionary state.  Denote by $N_{\rm phys}^{\rm m}(\delta;B)$ and
$N_{\rm phys}^{\rm p}(\delta;B)$ the least mixed- and counted-pure-start
label counts that approximate, in total variation to error at most
$\delta$, every common visible stopping and terminal-POVM policy whose
target expected calls are at most $B$.

\begin{theorem}[Sharp leading order for a fixed badly approximable phase]
\label{thm:r76-physical-logarithm}
Suppose $\omega$ is irrational and has a supplied constant $c>0$ such that
\begin{equation}\label{eq:r76-phase-separation}
 \|q\omega\|_{\mathbb R/\mathbb Z}\ge c/q
 \qquad(q\in\mathbb Z_{\ge1}).
\end{equation}
Necessarily $c\le1/2$.  Fix $B>0$ and put $\zeta=\min\{1,7B/8\}$.
For every integer $J\ge2$, if
\begin{equation}\label{eq:r76-physical-lower-condition}
 0<\delta\le\frac{\zeta}{8^J(1+J^2/c^2)},
\end{equation}
then
\begin{equation}\label{eq:r76-physical-lower}
 N_{\rm phys}^{\rm m}(\delta;B)\ge J,
 \qquad N_{\rm phys}^{\rm p}(\delta;B)\ge J+1.
\end{equation}

There is also a finite integer $K_\omega$, independent of $B,\delta$,
such that, with $L$ from \eqref{eq:r76-qubit-depth},
\begin{equation}\label{eq:r76-physical-upper}
 N_{\rm phys}^{\rm m}(\delta;B)\le K_\omega+L-1,
 \qquad N_{\rm phys}^{\rm p}(\delta;B)\le K_\omega+L.
\end{equation}
Consequently, for either start convention and every fixed $B>0$,
\begin{equation}\label{eq:r76-physical-log-limit}
 \lim_{\delta\downarrow0}
 \frac{N_{\rm phys}(\delta;B)}{\log_8(1/\delta)}=1.
\end{equation}
For $\omega=(\sqrt5-1)/2$, one may take $c=1/3$ and $K_\omega=5$.
All approximation lower bounds in this theorem are for physical
dictionaries with the stated Born terminal decoders.
\end{theorem}
\begin{proof}
For $n\ge0$, consider the report word $1\,0^n$ and the pure target state
$\rho_n$ it prepares.  Its Bloch vector is
$v_n=(\cos(2\pi n\omega),\sin(2\pi n\omega),0)$, with the coordinates
ordered along the real Bloch plane.  The word has probability
$8^{-(n+1)}$.  Use the following policy.  With probability $\zeta$ begin
the test; otherwise stop before any call.  During the test stop as soon
as a report differs from the specified prefix, or at its completion.
Apply the terminal POVM $(\rho_n,I-\rho_n)$ on successful completion.
The target expected number of calls is
\[
 \zeta\sum_{j=0}^{n}8^{-j}\le\zeta\,8/7\le B.
\]
Thus this is one of the policies in the definition of the approximation
error, for every $n$, without a horizon-dependent loss of the budget.

In the simulator let $A_n=\sum_i a_{n,i}\sigma_i$, $a_{n,i}\ge0$, be
the subnormalized terminal state on the event that this test is begun
and its complete prefix is observed.  Put
$p_n=\zeta 8^{-(n+1)}$.  Two events in the full stopped record give
\begin{equation}\label{eq:r76-prefix-mass-fidelity}
 |\tr A_n-p_n|\le\delta,
 \qquad \tr[(I-\rho_n)A_n]\le\delta.
\end{equation}
The second event has zero target probability.  If $p_n>\delta$, these
inequalities imply that some dictionary state with positive weight obeys
\[
 \tr(\rho_n\sigma_i)
 \ge 1-\frac{\delta}{p_n-\delta}.
\]
This conclusion is derived from joint events.  No error bound for a
rare conditional distribution has been assumed.

For $0\le n<J$, condition \eqref{eq:r76-physical-lower-condition} gives
\[
 \frac{\delta}{p_n-\delta}\le c^2/J^2.
\]
If $r_i\in\mathbb R^3$, $|r_i|\le1$, is the Bloch vector of the
corresponding physical dictionary state, then
\[
 |r_i-v_n|^2
 \le 2-2r_i\mathbin{\cdot}v_n
 =4[1-\tr(\rho_n\sigma_i)]\le4c^2/J^2.
\]
On the other hand, for $0\le m<n<J$, the elementary inequality
$\sin(\pi x)\ge2x$ on $[0,1/2]$ and
\eqref{eq:r76-phase-separation} give
\[
 |v_n-v_m|=2|\sin(\pi(n-m)\omega)|
 \ge4\|(n-m)\omega\|_{\mathbb R/\mathbb Z}
 \ge\frac{4c}{n-m}>\frac{4c}{J}.
\]
One dictionary vector cannot be within distance $2c/J$ of two such
points.  At least $J$ distinct labels are therefore necessary.

For a pure start, the initial label has its own physical density matrix
$\sigma_o$.  Stopping before any call and choosing any binary terminal
POVM is allowed at zero cost.  The error condition therefore implies
$\tfrac12\|\sigma_o-I/2\|_1\le\delta$, or $|r_o|\le2\delta$.
Its fidelity with any pure state is at most $1/2+\delta$.
Under \eqref{eq:r76-physical-lower-condition}, $J\ge2$ and $c\le1/2$
give $\delta<1/4$ and $c^2/J^2\le1/16$.  Thus the initial label cannot
be one of the $J$ labels whose fidelity is at least $15/16$.
This proves both inequalities in \eqref{eq:r76-physical-lower}.

For the improved upper, choose $K=K_\omega$ such that
\[
 P_K=\operatorname{conv}\{v_0,\ldots,v_{K-1}\}
 \supseteq\{x\in\mathbb R^2:|x|\le5/11\}.
\]
Such a finite $K$ exists for every irrational $\omega$.  Here is a direct
argument sufficient for this use.  Pigeonholing finitely many fractional
multiples of $\omega$ gives an integer $q>0$ with
$0<d=\|q\omega\|<1/4$.  The points with indices
$0,q,2q,\ldots,\lfloor1/d\rfloor q$ have successive circular gaps at
most $d$, in one orientation or the other.  Taking a prefix containing
all these indices, its inscribed polygon has every central gap less
than $\pi/2$, and hence inradius greater than $1/\sqrt2>5/11$.
The origin is in this polygon.  It is a valid seed polytope: report $1$
returns to $v_0$, and every report-$2$ output is in its contained disk.

For one angle, all continuation images of these $K$ seeds through depth
$L-1$ belong to
$\{v_0,\ldots,v_{K+L-2}\}$, a set of $K+L-1$ states.  Returning
branches are expanded in the original $K$ seeds.  No overflow can occur
before $L$ successive zero reports, even when a return chooses the last
seed.  Theorem~\ref{thm:r76-return-approximation} therefore proves
\eqref{eq:r76-physical-upper}.  The finite barycentric coefficients can
be chosen from any triangulation of $P_K$.  All labels, including the
optional initializer representing $I/2$, are physical states.

For fixed $B$, the upper has
$L=\log_8(1/\delta)+O_B(1)$.  Conversely,
\eqref{eq:r76-physical-lower-condition} holds for an integer
\[
 J=\log_8(1/\delta)
       -2\log_8\log_8(1/\delta)+O_{B,c}(1)
\]
chosen sufficiently small, as follows by taking logarithms of that
inequality.  Dividing the two bounds by $\log_8(1/\delta)$ proves
\eqref{eq:r76-physical-log-limit}.

Finally put $\omega=(\sqrt5-1)/2$, and let $p$ be a nearest integer to
$q\omega$.  Its conjugate is $\omega'=-(\sqrt5+1)/2$.  The product
$(q\omega-p)(q\omega'-p)=p^2+pq-q^2$ is a nonzero integer.  Therefore
\[
 |q\omega-p|\ge\frac1{|q\omega'-p|}
 \ge\frac1{\sqrt5q+1/2}>\frac1{3q},
\]
proving \eqref{eq:r76-phase-separation} with $c=1/3$.
The first five fractional orbit angles, in circular order, are
\[
 0,\quad 2\omega-1,\quad4\omega-2,\quad\omega,\quad3\omega-1.
\]
Their successive gaps are $2\omega-1$ or $2-3\omega$.
Since $3/5<\omega<5/8$, their maximum is $2\omega-1<1/4$.
This proves the required polygon inclusion with $K_\omega=5$.
\end{proof}

The final limit fixes $B>0$; it is not a uniform two-parameter asymptotic
as $B$ also varies.  For arbitrary classical decoder tables, the general
finite-alphabet upper remains valid, while the present packing lower
need not apply.  The badly approximable phase condition has an explicit
quadratic-irrational instance above.  It is not asserted for the separate
algebraic rotation $(3+4i)/5$ studied earlier in the paper.

The improved upper does apply to that original algebraic rotation, with
an explicit constant.  Write $\theta=\arg((3+4i)/5)\in(0,\pi/2)$.
The inequalities $\pi/6<\theta<\pi/3$ put the first seven orbit points
$0,\theta,\ldots,6\theta$ in their circular order with every gap less
than $\pi$.  The six consecutive edges have distance
$\cos(\theta/2)=2/\sqrt5$ from the origin.  The closing edge has
distance $-\cos(3\theta)=117/125$.  Both exceed $5/11$, so $K=7$ is
valid.  This gives $L+6$ mixed labels or $L+7$ pure labels at the error
in \eqref{eq:r76-qubit-error}.  All those orbit coordinates are rational.
Triangulating their rational polygon gives rational initial and returning
barycentric rows, so the resulting finite transition tables can be chosen
rational.  This is an explicit upper for that algebraic instrument, not
a claim about its unresolved Diophantine lower certificate.
